\documentclass[11pt,a4paper]{article}

% ============================================================
%  QIMMA -- COURS : Structures algébriques
%  2ème Bac Sciences Mathématiques (A et B)
%  Standard exhaustif : définitions formelles + démonstrations
%  Basé sur templates/qimma-template.tex (charte Qimma)
% ============================================================

\usepackage[french]{babel}
\usepackage[utf8]{inputenc}
\usepackage[T1]{fontenc}
\usepackage{lmodern}
\usepackage[a4paper,margin=2cm,top=2.3cm,bottom=2.6cm]{geometry}
\usepackage{amsmath,amssymb}
\usepackage{xcolor}
\usepackage{graphicx}
\usepackage{tikz}
\usetikzlibrary{shadows,calc}
\usepackage[most]{tcolorbox}
\usepackage{titlesec}
\usepackage{fancyhdr}
\usepackage{enumitem}
\usepackage{pifont}
\usepackage{array}
\usepackage{colortbl}
\usepackage{hyperref}

% ------------------- CHARTE QIMMA -------------------
\definecolor{qteal}{HTML}{0d9488}
\definecolor{qtealdark}{HTML}{134e4a}
\definecolor{qamber}{HTML}{fbbf24}
\definecolor{qambertext}{HTML}{92400e}
\definecolor{ink}{HTML}{1f2937}
\definecolor{inkgray}{HTML}{6b7280}
\definecolor{tealpale}{HTML}{e6f5f3}
\colorlet{colDef}{qteal}
\definecolor{colProp}{HTML}{0f766e}
\definecolor{colEx}{HTML}{b45309}
\definecolor{colRem}{HTML}{9d174d}
\color{ink}

\graphicspath{{../../../../assets/}}
\newcommand{\logofile}{../../../../assets/qimma-logo.pdf}

% ------------------- METADONNEES -------------------
\newcommand{\matiere}{Mathématiques}
\newcommand{\niveau}{2\up{ème} Bac -- Sciences Mathématiques}
\newcommand{\chapitretitre}{Structures algébriques}

% ------------------- EN-TETE / PIED DE PAGE -------------------
\pagestyle{fancy}
\fancyhf{}
\renewcommand{\headrulewidth}{0pt}
\renewcommand{\footrulewidth}{0.5pt}
\fancyfoot[L]{\raisebox{-0.15cm}{\includegraphics[height=0.35cm]{\logofile}}~\small\sffamily\color{inkgray}Qimma}
\fancyfoot[C]{\small\sffamily\color{inkgray}\matiere\ -- \niveau}
\fancyfoot[R]{\small\sffamily\color{qtealdark}\thepage}

% ------------------- TITRES DE SECTION -------------------
\titleformat{\section}
  {\normalfont\sffamily\Large\bfseries\color{qtealdark}}
  {\tikz[baseline=-3.2pt]{\node[circle,fill=qteal,text=white,inner sep=0pt,minimum size=8mm]{\sffamily\bfseries\thesection};}}
  {10pt}{}
  [\vspace{-2mm}{\color{qamber}\titlerule[1.6pt]}\vspace{2mm}]
\titlespacing*{\section}{0pt}{16pt}{10pt}

\titleformat{\subsection}
  {\normalfont\sffamily\large\bfseries\color{qtealdark}}
  {\color{qteal}\ding{212}}{6pt}{}
\titlespacing*{\subsection}{0pt}{12pt}{6pt}

% ------------------- BOITES PEDAGOGIQUES -------------------
\newtcolorbox{fiche}[3][]{
  enhanced, breakable,
  colback=white, colframe=#2,
  boxrule=1pt, arc=2mm,
  top=7mm, left=4mm, right=4mm, bottom=3mm,
  drop shadow={black!25, shadow xshift=1mm, shadow yshift=-1mm},
  attach boxed title to top left={xshift=4mm, yshift=-3mm, yshifttext=-1mm},
  boxed title style={colback=#2, arc=1.5mm, boxrule=0pt, left=2mm, right=2mm},
  coltitle=white, fonttitle=\sffamily\bfseries,
  title={#3}, #1
}
\newenvironment{definition}[1][Définition]
  {\begin{fiche}{colDef}{\ding{110}~~#1}}{\end{fiche}}
\newenvironment{propriete}[1][Propriété]
  {\begin{fiche}{colProp}{\ding{72}~~#1}}{\end{fiche}}
\newenvironment{exemple}[1][Exemple]
  {\begin{fiche}{colEx}{\ding{217}~~#1}}{\end{fiche}}
\newenvironment{remarque}[1][Remarque]
  {\begin{fiche}{colRem}{\ding{43}~~#1}}{\end{fiche}}

% Démonstration (hors boîte, style discret)
\newenvironment{demo}
  {\par\smallskip\noindent{\sffamily\bfseries\color{qtealdark}Démonstration.}\ \itshape}
  {\ \normalfont\hfill$\blacksquare$\par\medskip}

\hypersetup{colorlinks=true, linkcolor=qtealdark, urlcolor=qtealdark}

% ============================================================
\begin{document}

% ------------------- BANNIERE DE TITRE -------------------
\begin{tcolorbox}[
  enhanced, boxrule=0pt, arc=4mm,
  interior style={left color=qtealdark, right color=qteal},
  top=8mm, bottom=8mm, left=10mm, right=32mm,
  drop shadow={black!35, shadow xshift=1.5mm, shadow yshift=-1.5mm},
  before skip=0pt, after skip=9mm,
  overlay={
    \node[anchor=north west] at ([xshift=6mm,yshift=-6mm]frame.north west)
      {\includegraphics[height=1.25cm]{\logofile}};
    \node[anchor=north east, fill=qamber, text=qambertext, rounded corners=2pt,
          inner sep=4pt, font=\sffamily\bfseries\small] at ([xshift=-6mm,yshift=-6mm]frame.north east) {COURS};
  }
]
\hspace{1.6cm}\centering
{\sffamily\bfseries\color{white}\fontsize{23}{27}\selectfont \chapitretitre}\\[3mm]
{\sffamily\color{white!90}\large \matiere\ \textbullet\ \niveau}
\end{tcolorbox}

% ------------------- OBJECTIFS -------------------
\begin{tcolorbox}[
  enhanced, colback=tealpale, colframe=qteal, boxrule=0.8pt, arc=2mm,
  left=5mm, right=5mm, top=3mm, bottom=3mm,
  title={\sffamily\bfseries\color{qtealdark}\ding{51}~~Objectifs du chapitre},
  coltitle=qtealdark, colbacktitle=tealpale, boxrule=0.8pt,
  attach title to upper={\par\vspace{1mm}}
]
\begin{itemize}[leftmargin=6mm, label=\ding{212}, itemsep=1mm]
  \item Définir une loi de composition interne et étudier ses propriétés : associativité, commutativité, élément neutre, symétrique.
  \item Reconnaître et démontrer qu'un ensemble muni d'une loi est un groupe, et manipuler les groupes usuels.
  \item Étudier la structure d'anneau et de corps, et connaître les exemples fondamentaux ($\mathbb{Z}$, $\mathbb{Q}$, $\mathbb{R}$, $\mathbb{C}$, $\mathbb{Z}/n\mathbb{Z}$).
  \item Démontrer que $(\mathbb{Z}/n\mathbb{Z},+,\times)$ est un anneau, et un corps si et seulement si $n$ est premier.
  \item Démontrer les propriétés d'un sous-groupe à l'aide de caractérisations pratiques.
  \item Définir un morphisme de groupes, démontrer ses propriétés élémentaires, et caractériser son injectivité par le noyau.
  \item Résoudre les exercices type examen national : tables de lois, vérification d'axiomes, morphismes simples.
\end{itemize}
\end{tcolorbox}

\vspace{4mm}

% ------------------- SECTION 1 -------------------
\section{Lois de composition interne}

\begin{definition}[Loi de composition interne]
Soit $E$ un ensemble non vide. Une \textbf{loi de composition interne} sur $E$ est une application $*:E\times E\to E$, qui à $(x,y)$ associe $x*y\in E$. « Interne » signifie que le résultat reste dans $E$.
\end{definition}

\begin{definition}[Propriétés d'une loi]
Soit $*$ une loi sur $E$.
\begin{itemize}[leftmargin=6mm, itemsep=0.5mm]
  \item $*$ est \textbf{associative} si $(\forall x,y,z\in E)\ (x*y)*z=x*(y*z)$.
  \item $*$ est \textbf{commutative} si $(\forall x,y\in E)\ x*y=y*x$.
  \item $e\in E$ est \textbf{élément neutre} pour $*$ si $(\forall x\in E)\ x*e=e*x=x$.
  \item Si $*$ admet un neutre $e$, $x'\in E$ est \textbf{symétrique} de $x$ si $x*x'=x'*x=e$.
\end{itemize}
\end{definition}

\begin{propriete}[Unicité du neutre et du symétrique]
\begin{itemize}[leftmargin=6mm, itemsep=0.5mm]
  \item Si $*$ admet un élément neutre, il est \textbf{unique}.
  \item Si $*$ est associative et admet un neutre, le symétrique d'un élément, s'il existe, est \textbf{unique}.
\end{itemize}
\end{propriete}

\begin{demo}
\emph{Unicité du neutre} : soient $e,e'$ deux neutres. Alors $e*e'=e$ (car $e'$ neutre) et $e*e'=e'$ (car $e$ neutre) : donc $e=e'$.

\smallskip
\emph{Unicité du symétrique} : soient $x',x''$ deux symétriques de $x$. Alors
\[
  x' = x'*e = x'*(x*x'') = (x'*x)*x'' = e*x'' = x''
\]
(en utilisant l'associativité et les définitions de $x'$, $x''$ comme symétriques).
\end{demo}

\begin{exemple}[Étude d'une loi non usuelle]
Sur $E=\mathbb{R}$, on définit $x*y=x+y-1$. Étudions cette loi.

\smallskip
\emph{Interne} : $x*y\in\mathbb{R}$ pour tous $x,y\in\mathbb{R}$. $\checkmark$

\smallskip
\emph{Commutative} : $x*y=x+y-1=y+x-1=y*x$. $\checkmark$

\smallskip
\emph{Associative} : $(x*y)*z=(x+y-1)*z=(x+y-1)+z-1=x+y+z-2$ ; $x*(y*z)=x+(y+z-1)-1=x+y+z-2$. Égaux. $\checkmark$

\smallskip
\emph{Neutre} : cherchons $e$ tel que $x*e=x$ pour tout $x$ : $x+e-1=x\iff e=1$. Le neutre est $e=1$.

\smallskip
\emph{Symétrique} : pour $x$ fixé, cherchons $x'$ tel que $x*x'=1$ : $x+x'-1=1\iff x'=2-x$. Tout élément admet un symétrique.
\end{exemple}

% ------------------- SECTION 2 -------------------
\section{Structure de groupe}

\begin{definition}[Groupe]
$(G,*)$ est un \textbf{groupe} lorsque :
\begin{enumerate}[leftmargin=7mm, itemsep=0.5mm]
  \item $*$ est une loi de composition interne sur $G$ ;
  \item $*$ est \textbf{associative} ;
  \item $*$ admet un \textbf{élément neutre} $e$ dans $G$ ;
  \item tout élément de $G$ admet un \textbf{symétrique} dans $G$ pour $*$.
\end{enumerate}
Si de plus $*$ est \textbf{commutative}, $(G,*)$ est un \textbf{groupe commutatif} (ou abélien).
\end{definition}

\begin{exemple}[Groupes usuels]
\begin{itemize}[leftmargin=6mm, itemsep=0.5mm]
  \item $(\mathbb{Z},+)$, $(\mathbb{Q},+)$, $(\mathbb{R},+)$, $(\mathbb{C},+)$ sont des groupes commutatifs (neutre $0$, symétrique de $x$ : $-x$).
  \item $(\mathbb{Q}^*,\times)$, $(\mathbb{R}^*,\times)$, $(\mathbb{C}^*,\times)$ sont des groupes commutatifs (neutre $1$, symétrique de $x$ : $\frac1x$).
  \item $(\mathbb{N},+)$ \textbf{n'est pas} un groupe : les éléments non nuls n'ont pas de symétrique dans $\mathbb{N}$ (ex. $3$ n'a pas d'opposé dans $\mathbb{N}$).
  \item $(\mathbb{Z},\times)$ \textbf{n'est pas} un groupe : $2$ n'a pas d'inverse dans $\mathbb{Z}$ ($\frac12\notin\mathbb{Z}$).
\end{itemize}
\end{exemple}

\begin{propriete}[Simplification dans un groupe]
Dans un groupe $(G,*)$, pour $a,b,c\in G$ :
\[
  a*b=a*c \implies b=c \qquad \text{(et de même } b*a=c*a\implies b=c\text{)}.
\]
\end{propriete}

\begin{demo}
En composant à gauche par le symétrique $a'$ de $a$ : $a'*(a*b)=a'*(a*c)$, soit par associativité $(a'*a)*b=(a'*a)*c$, c'est-à-dire $e*b=e*c$, donc $b=c$.
\end{demo}

\begin{definition}[Sous-groupe]
Soit $(G,*)$ un groupe et $H\subset G$, $H\ne\varnothing$. $H$ est un \textbf{sous-groupe} de $G$ lorsque $(H,*)$ est lui-même un groupe.
\end{definition}

\begin{propriete}[Caractérisation pratique d'un sous-groupe]
$H\subset G$, $H\ne\varnothing$, est un sous-groupe de $(G,*)$ si et seulement si :
\[
  (\forall x,y\in H)\quad x*y'\in H
\]
où $y'$ désigne le symétrique de $y$ dans $G$ (formulation additive : $x-y\in H$).
\end{propriete}

\begin{exemple}[Exemple de sous-groupe]
Montrons que $2\mathbb{Z}=\{2k : k\in\mathbb{Z}\}$ est un sous-groupe de $(\mathbb{Z},+)$.

\smallskip
$2\mathbb{Z}\ne\varnothing$ ($0\in2\mathbb{Z}$). Soient $x=2k$, $y=2k'\in2\mathbb{Z}$ : $x-y=2(k-k')\in2\mathbb{Z}$. Donc $2\mathbb{Z}$ est un sous-groupe de $(\mathbb{Z},+)$.
\end{exemple}

% ------------------- SECTION 3 -------------------
\section{Morphismes de groupes}

\begin{definition}[Morphisme de groupes]
Soient $(G,*)$ et $(G',\top)$ deux groupes. Une application $f:G\to G'$ est un \textbf{morphisme de groupes} (ou homomorphisme) lorsque
\[
  (\forall x,y\in G)\quad f(x*y) = f(x)\top f(y).
\]
Un morphisme \textbf{bijectif} est appelé \textbf{isomorphisme} ; un morphisme de $G$ dans lui-même est un \textbf{endomorphisme}.
\end{definition}

\begin{propriete}[Propriétés élémentaires]
Soit $f:(G,*)\to(G',\top)$ un morphisme de groupes, de neutres respectifs $e$ et $e'$. Alors :
\[
  f(e) = e', \qquad (\forall x\in G)\ f(x^{-1}) = \bigl(f(x)\bigr)^{-1},
\]
où $x^{-1}$ désigne le symétrique de $x$ dans $(G,*)$.
\end{propriete}

\begin{demo}
$f(e)=f(e*e)=f(e)\top f(e)$. En composant à gauche par le symétrique de $f(e)$ dans $G'$ (règle de simplification dans un groupe, vue en section 2) : $e'=f(e)$.

\smallskip
Ensuite, $f(x)\top f(x^{-1}) = f(x*x^{-1}) = f(e) = e'$, donc $f(x^{-1})$ est bien le symétrique de $f(x)$ dans $G'$, c'est-à-dire $f(x^{-1})=\bigl(f(x)\bigr)^{-1}$.
\end{demo}

\begin{definition}[Noyau et image]
Le \textbf{noyau} de $f$ est $\ker f=\{x\in G : f(x)=e'\}$ ; l'\textbf{image} de $f$ est $\operatorname{Im}f=\{f(x) : x\in G\}$.
\end{definition}

\begin{propriete}[Caractérisation de l'injectivité, sous-groupes associés]
$\ker f$ est un sous-groupe de $(G,*)$, $\operatorname{Im}f$ est un sous-groupe de $(G',\top)$, et
\[
  f \text{ est injectif} \iff \ker f=\{e\}.
\]
\end{propriete}

\begin{demo}
($\Rightarrow$) Si $f$ injectif et $x\in\ker f$ : $f(x)=e'=f(e)$, donc $x=e$ par injectivité : $\ker f=\{e\}$.

\smallskip
($\Leftarrow$) Si $\ker f=\{e\}$ et $f(x)=f(y)$ : alors $f(x)\top\bigl(f(y)\bigr)^{-1}=e'$, soit $f(x)\top f(y^{-1})=e'$, c'est-à-dire $f(x*y^{-1})=e'$ : donc $x*y^{-1}\in\ker f=\{e\}$, soit $x*y^{-1}=e$, d'où $x=y$.
\end{demo}

\begin{exemple}[Morphisme classique]
L'application $\exp:(\mathbb{R},+)\to(\mathbb{R}_+^*,\times)$, $x\mapsto e^x$, est un morphisme de groupes : $\exp(x+y)=e^{x+y}=e^xe^y=\exp(x)\times\exp(y)$. C'est même un \textbf{isomorphisme} (bijection de $\mathbb{R}$ sur $\mathbb{R}_+^*$, vue au chapitre sur la fonction exponentielle) : $\ker(\exp)=\{0\}$, cohérent avec l'injectivité de $\exp$.
\end{exemple}

\begin{exemple}[Le module comme morphisme]
L'application $|\cdot|:(\mathbb{C}^*,\times)\to(\mathbb{R}_+^*,\times)$, $z\mapsto|z|$, est un morphisme de groupes : $|zz'|=|z||z'|$ (propriété démontrée au chapitre sur les nombres complexes). Ce morphisme n'est \textbf{pas} injectif : $\ker(|\cdot|)=\{z\in\mathbb{C}^* : |z|=1\}=\mathbb{U}$, le sous-groupe des complexes de module $1$ (déjà rencontré comme exemple de sous-groupe de $(\mathbb{C}^*,\times)$).
\end{exemple}

% ------------------- SECTION 4 -------------------
\section{Structures d'anneau et de corps}

\begin{definition}[Anneau]
$(A,+,\times)$ est un \textbf{anneau} lorsque :
\begin{enumerate}[leftmargin=7mm, itemsep=0.5mm]
  \item $(A,+)$ est un \textbf{groupe commutatif} (neutre noté $0_A$) ;
  \item $\times$ est \textbf{associative} et possède un élément \textbf{neutre} $1_A$ ;
  \item $\times$ est \textbf{distributive} par rapport à $+$ : $x\times(y+z)=x\times y+x\times z$ et $(y+z)\times x=y\times x+z\times x$.
\end{enumerate}
Si de plus $\times$ est commutative, $A$ est un \textbf{anneau commutatif}.
\end{definition}

\begin{exemple}[Anneaux usuels]
$(\mathbb{Z},+,\times)$, $(\mathbb{Q},+,\times)$, $(\mathbb{R},+,\times)$, $(\mathbb{C},+,\times)$ sont des anneaux commutatifs.
\end{exemple}

\begin{definition}[Corps]
$(K,+,\times)$ est un \textbf{corps} lorsque $(K,+,\times)$ est un anneau \textbf{commutatif}, avec $0_K\ne1_K$, et tel que tout élément \textbf{non nul} admette un symétrique pour $\times$ (c'est-à-dire $(K^*,\times)$ est un groupe).
\end{definition}

\begin{exemple}[Corps et non-corps]
$(\mathbb{Q},+,\times)$, $(\mathbb{R},+,\times)$, $(\mathbb{C},+,\times)$ sont des \textbf{corps}. En revanche $(\mathbb{Z},+,\times)$ \textbf{n'est pas} un corps : $2\in\mathbb{Z}^*$ n'a pas d'inverse dans $\mathbb{Z}$ pour $\times$.
\end{exemple}

\begin{propriete}[La structure $(\mathbb{Z}/n\mathbb{Z},+,\times)$ -- lien avec l'arithmétique]
Munissons $\mathbb{Z}/n\mathbb{Z}=\{\overline0,\overline1,\ldots,\overline{n-1}\}$ (les classes de congruence modulo $n$, chapitre précédent) des opérations $\overline a+\overline b=\overline{a+b}$ et $\overline a\times\overline b=\overline{a\times b}$ (bien définies grâce à la compatibilité des congruences avec $+$ et $\times$). Alors :
\begin{itemize}[leftmargin=6mm, itemsep=0.5mm]
  \item $(\mathbb{Z}/n\mathbb{Z},+,\times)$ est toujours un \textbf{anneau commutatif}, quel que soit $n\geq2$.
  \item C'est un \textbf{corps} si et seulement si $n$ est \textbf{premier}.
\end{itemize}
\end{propriete}

\begin{demo}
\emph{Anneau} : $(\mathbb{Z}/n\mathbb{Z},+)$ est un groupe commutatif (neutre $\overline0$, symétrique de $\overline a$ : $\overline{-a}$) ; $\times$ est associative, commutative, de neutre $\overline1$, et distributive par rapport à $+$ (ces propriétés se transmettent directement de $(\mathbb{Z},+,\times)$ via le passage aux classes).

\smallskip
\emph{Corps si $n$ premier} : si $n=p$ premier, soit $\overline a\ne\overline0$, c'est-à-dire $p\nmid a$, donc $a\wedge p=1$ (car $p$ premier). Par Bézout, il existe $u,v$ tels que $au+pv=1$, soit $au\equiv1\,[p]$ : $\overline u$ est l'inverse de $\overline a$. Donc tout élément non nul est inversible : $\mathbb{Z}/p\mathbb{Z}$ est un corps.

\smallskip
\emph{Pas un corps si $n$ n'est pas premier} : si $n=ab$ avec $1<a,b<n$, alors $\overline a\times\overline b=\overline{ab}=\overline n=\overline0$, avec $\overline a\ne\overline0$ et $\overline b\ne\overline0$ : $\overline a$ n'est pas inversible (un élément inversible ne peut pas être facteur de $\overline0$ avec un facteur non nul), donc $\mathbb{Z}/n\mathbb{Z}$ n'est pas un corps.
\end{demo}

\begin{exemple}[Calcul dans $\mathbb{Z}/6\mathbb{Z}$ et $\mathbb{Z}/5\mathbb{Z}$]
Dans $\mathbb{Z}/6\mathbb{Z}$ : $\overline2\times\overline3=\overline6=\overline0$, alors que $\overline2\ne\overline0$ et $\overline3\ne\overline0$ : $\mathbb{Z}/6\mathbb{Z}$ n'est pas un corps ($6=2\times3$ n'est pas premier).

\smallskip
Dans $\mathbb{Z}/5\mathbb{Z}$ (5 premier) : cherchons l'inverse de $\overline3$. Par Bézout, $3\times2-5\times1=1$, donc $\overline3\times\overline2=\overline1$ : l'inverse de $\overline3$ est $\overline2$.
\end{exemple}

\begin{remarque}[Hiérarchie des structures]
Un corps est un anneau (commutatif, avec en plus l'inversibilité des éléments non nuls) ; un anneau contient un groupe commutatif $(A,+)$. La hiérarchie est : \textbf{groupe} $\subset$ \textbf{anneau} $\subset$ \textbf{corps} (au sens où chaque structure ajoute des axiomes).
\end{remarque}

% ------------------- SECTION 5 -------------------
\section{Exercices corrigés -- niveau examen national SM}

\begin{exemple}[Exercice 1 -- vérification complète d'une structure de groupe]
Sur $\mathbb{R}$, on définit $x*y=x+y+xy$. Montrer que $(\mathbb{R}\setminus\{-1\},*)$ est un groupe commutatif.

\smallskip
\textbf{Corrigé.} \emph{Interne} : pour $x,y\ne-1$, montrons $x*y\ne-1$. Supposons $x*y=-1$ : $x+y+xy=-1\iff(x+1)(y+1)=0$ (développement : $xy+x+y+1=(x+1)(y+1)$), ce qui contredit $x,y\ne-1$. Donc $x*y\ne-1$ : la loi est bien interne sur $\mathbb{R}\setminus\{-1\}$.

\smallskip
\emph{Commutative} : $x*y=x+y+xy=y+x+yx=y*x$. $\checkmark$

\smallskip
\emph{Associative} : $(x*y)*z=(x+y+xy)+z+(x+y+xy)z=x+y+z+xy+xz+yz+xyz$ ; par symétrie du calcul, $x*(y*z)$ donne la même expression. $\checkmark$ (admis par symétrie de l'expression développée).

\smallskip
\emph{Neutre} : $x*e=x\iff x+e+xe=x\iff e(1+x)=0$. Comme $x\ne-1$, on a $e=0$. Vérification : $0\in\mathbb{R}\setminus\{-1\}$. $\checkmark$

\smallskip
\emph{Symétrique} : $x*x'=0\iff x+x'+xx'=0\iff x'(1+x)=-x\iff x'=\dfrac{-x}{1+x}$ (licite car $x\ne-1$). Vérifions $x'\ne-1$ : $\frac{-x}{1+x}=-1\iff-x=-1-x\iff0=-1$, faux : donc $x'\ne-1$ toujours.

\smallskip
\textbf{Conclusion} : $(\mathbb{R}\setminus\{-1\},*)$ est un groupe commutatif, de neutre $0$ et de symétrique $x'=\dfrac{-x}{1+x}$.
\end{exemple}

\begin{exemple}[Exercice 2 -- sous-groupe des racines de l'unité]
Montrer que l'ensemble $\mathbb{U}=\{z\in\mathbb{C} : |z|=1\}$ est un sous-groupe de $(\mathbb{C}^*,\times)$.

\smallskip
\textbf{Corrigé.} $\mathbb{U}\ne\varnothing$ ($1\in\mathbb{U}$). Soient $z,z'\in\mathbb{U}$ : montrons $z\times\dfrac1{z'}\in\mathbb{U}$.
\[
  \left|\frac{z}{z'}\right| = \frac{|z|}{|z'|} = \frac11 = 1.
\]
Donc $\dfrac{z}{z'}\in\mathbb{U}$ : $\mathbb{U}$ est un sous-groupe de $(\mathbb{C}^*,\times)$.
\end{exemple}

\begin{exemple}[Exercice 3 -- table de loi sur un ensemble fini]
Sur $E=\{1,-1,i,-i\}\subset\mathbb{C}$, on considère la multiplication usuelle des complexes. Dresser la table et vérifier que $(E,\times)$ est un groupe.

\smallskip
\textbf{Corrigé.}
\begin{center}
\renewcommand{\arraystretch}{1.3}
\begin{tabular}{c|cccc}
$\times$ & $1$ & $-1$ & $i$ & $-i$ \\ \hline
$1$ & $1$ & $-1$ & $i$ & $-i$ \\
$-1$ & $-1$ & $1$ & $-i$ & $i$ \\
$i$ & $i$ & $-i$ & $-1$ & $1$ \\
$-i$ & $-i$ & $i$ & $1$ & $-1$ \\
\end{tabular}
\end{center}
La table montre que le produit de deux éléments de $E$ reste dans $E$ (loi interne). Le neutre est $1$ (colonne/ligne identité). Chaque ligne contient $1$ : chaque élément a un symétrique ($1$ et $-1$ sont leur propre symétrique, $i$ et $-i$ sont symétriques l'un de l'autre). L'associativité vient de celle de $\times$ dans $\mathbb{C}$. Donc $(E,\times)$ est un groupe (commutatif, car la table est symétrique par rapport à la diagonale).
\end{exemple}

\begin{exemple}[Exercice 4 -- morphisme, noyau et injectivité]
On munit $\mathbb{R}$ de $+$ et $\mathbb{U}=\{z\in\mathbb{C}:|z|=1\}$ de $\times$ (groupe, cf. exercice 2). Soit $f:(\mathbb{R},+)\to(\mathbb{U},\times)$ définie par $f(\theta)=e^{i\theta}$. Montrer que $f$ est un morphisme de groupes, déterminer $\ker f$, et dire si $f$ est injectif.

\smallskip
\textbf{Corrigé.} \emph{Morphisme} : pour $\theta,\theta'\in\mathbb{R}$, $f(\theta+\theta')=e^{i(\theta+\theta')}=e^{i\theta}e^{i\theta'}=f(\theta)\times f(\theta')$. $\checkmark$

\smallskip
\emph{Noyau} : $\ker f=\{\theta\in\mathbb{R} : e^{i\theta}=1\}=\{\theta\in\mathbb{R} : \theta\equiv0\ [2\pi]\}=2\pi\mathbb{Z}$.

\smallskip
\emph{Injectivité} : $\ker f\ne\{0\}$ (par exemple $2\pi\in\ker f$), donc d'après la caractérisation de l'injectivité, $f$ \textbf{n'est pas injectif} (cohérent : $f(0)=f(2\pi)=1$).
\end{exemple}

% ------------------- SECTION 6 -------------------
\section{Méthodes et pièges : la boîte à outils de l'examen}

\subsection{Ce qu'on te demande, ce que tu fais}

\begin{center}
\renewcommand{\arraystretch}{1.45}
\sffamily\small
\begin{tabular}{>{\raggedright\arraybackslash}p{7.2cm} >{\raggedright\arraybackslash}p{8cm}}
\rowcolor{qtealdark} \color{white}\bfseries Ce qu'on te demande & \color{white}\bfseries Ton réflexe \\
\rowcolor{tealpale} Montrer qu'une loi est interne & Vérifier que le résultat de l'opération reste dans l'ensemble de départ. \\
Trouver le neutre d'une loi & Résoudre $x*e=x$ (l'équation doit être vraie pour tout $x$, donc $e$ ne doit pas dépendre de $x$). \\
\rowcolor{tealpale} Trouver le symétrique de $x$ & Résoudre $x*x'=e$ d'inconnue $x'$. \\
Montrer que $(G,*)$ est un groupe & Vérifier dans l'ordre : loi interne, associativité, neutre, symétrique de chaque élément. \\
\rowcolor{tealpale} Montrer que $H$ est un sous-groupe de $G$ & $H\ne\varnothing$, puis $\forall x,y\in H,\ x*y'\in H$ (caractérisation pratique). \\
Montrer que $(K,+,\times)$ est un corps & Groupe commutatif pour $+$, anneau pour $\times$ avec neutre, tout élément non nul inversible. \\
\rowcolor{tealpale} Montrer que $f$ est un morphisme de groupes & Vérifier $f(x*y)=f(x)\top f(y)$ pour tous $x,y$. \\
Étudier l'injectivité de $f$ (morphisme) & Calculer $\ker f$ ; injectif $\iff\ker f=\{e\}$. \\
\rowcolor{tealpale} Montrer que $\mathbb{Z}/n\mathbb{Z}$ est un corps & Vérifier que $n$ est \textbf{premier} (sinon, chercher deux facteurs non nuls de produit $\overline0$). \\
\end{tabular}
\end{center}

\subsection{Les pièges qui coûtent des points}

\begin{remarque}[Erreurs relevées chaque année par les correcteurs]
\begin{itemize}[leftmargin=6mm, itemsep=0.5mm]
  \item Oublier de vérifier que la loi est \textbf{interne} avant toute autre propriété (souvent l'étape la plus délicate).
  \item Confondre \textbf{neutre} (élément qui laisse invariant) et \textbf{symétrique} (élément qui « annule » un autre).
  \item Oublier de vérifier que le neutre trouvé appartient bien à l'ensemble considéré (piège classique dans les exercices avec ensemble restreint, ex. $\mathbb{R}\setminus\{-1\}$).
  \item Se contenter de vérifier l'associativité sur un exemple numérique au lieu de la démontrer en toute généralité.
  \item Oublier qu'un anneau n'est \textbf{pas} nécessairement un corps ($\mathbb{Z}$ est un anneau, pas un corps).
  \item Pour un sous-groupe : oublier de vérifier $H\ne\varnothing$ avant d'appliquer la caractérisation pratique.
  \item Pour un morphisme : oublier de préciser les \textbf{deux} lois utilisées ($*$ dans $G$, $\top$ dans $G'$), ou confondre $\ker f$ (sous-ensemble de $G$) et $\operatorname{Im}f$ (sous-ensemble de $G'$).
  \item Conclure à l'injectivité d'un morphisme sans avoir calculé $\ker f$ dans son intégralité.
\end{itemize}
\end{remarque}

% ------------------- RESUME -------------------
\vspace{4mm}
\begin{tcolorbox}[
  enhanced, boxrule=0pt, arc=2mm,
  interior style={left color=qtealdark, right color=qteal},
  left=6mm, right=6mm, top=4mm, bottom=4mm,
  fontupper=\color{white}\sffamily,
  title={\sffamily\bfseries\color{white}\ding{72}~~À retenir},
  colbacktitle=qtealdark, coltitle=white, boxrule=0pt
]
\begin{itemize}[leftmargin=6mm, label={\color{qamber}\ding{212}}, itemsep=1mm]
  \item Loi interne : associativité, commutativité, neutre (unique), symétrique (unique si associative).
  \item Groupe : loi interne + associative + neutre + tout élément symétrisable ; commutatif si en plus $x*y=y*x$.
  \item Sous-groupe : $H\ne\varnothing$ et $\forall x,y\in H,\ x*y'\in H$.
  \item Morphisme $f:(G,*)\to(G',\top)$ : $f(x*y)=f(x)\top f(y)$ ; alors $f(e)=e'$, $f(x^{-1})=f(x)^{-1}$ ; injectif $\iff\ker f=\{e\}$.
  \item Anneau : $(A,+)$ groupe commutatif, $\times$ associative avec neutre, distributivité.
  \item Corps : anneau commutatif où tout élément non nul est inversible pour $\times$.
  \item Exemples clés : $(\mathbb{Z},+)$, $(\mathbb{R}^*,\times)$ groupes ; $\mathbb{Z}$ anneau non corps ; $\mathbb{Q},\mathbb{R},\mathbb{C}$ corps.
  \item $(\mathbb{Z}/n\mathbb{Z},+,\times)$ : toujours anneau commutatif ; \textbf{corps} $\iff n$ \textbf{premier} (sinon diviseurs de zéro non inversibles).
\end{itemize}
\end{tcolorbox}

\end{document}
